\documentclass[11pt]{article}

\usepackage[margin=1in]{geometry}
\usepackage{amsmath, amssymb}
\usepackage{graphicx}
\usepackage{booktabs}
\usepackage{hyperref}
\usepackage{natbib}
\usepackage{microtype}
\usepackage{setspace}
\usepackage{titlesec}

\onehalfspacing

\title{Continuous Thought Machines for Streaming Reasoning\\in Financial Microstructure}
\author{Asray Gopa\\
  Department of Computer Science\\
  Iowa State University\\
  \texttt{asray@iastate.edu}}
\date{}

\begin{document}

\maketitle

\begin{abstract}
We investigate whether the Continuous Thought Machine (CTM) paradigm---where a
model performs a variable number of internal computational ticks per external
observation---offers advantages for high-frequency limit order book (LOB)
mid-price direction prediction. Standard deep learning approaches to this task
are episodic, segmenting the data stream into fixed windows, which is
misaligned with the fundamentally online nature of market microstructure data.
We formulate the prediction task as strictly streaming classification on
Bitcoin perpetual futures data sourced from Tardis.dev, resampled at 50\,ms
intervals with a 5-second prediction horizon. We compare two CTM-style
architectures against logistic regression, LSTM, and Transformer baselines,
and further evaluate all models under synthetic market shocks including spread
widening and liquidity withdrawal. Our results provide a controlled
characterization of the continuous-thought paradigm in a real-world sequential
prediction domain.
\end{abstract}

% ─────────────────────────────────────────────────────────────
\section{Introduction}
% ─────────────────────────────────────────────────────────────

High-frequency financial markets generate a continuous stream of information
in the form of limit order book (LOB) updates, where prices and quantities at
multiple depth levels evolve on millisecond timescales. Predicting short-horizon
mid-price direction from this stream---a canonical microstructure prediction
task---has attracted significant attention due to its theoretical and practical
importance~\citep{gould2013limit, cont2010stochastic}. The problem is naturally
sequential: each new order book event arrives as a single timestep, and a model
must update its belief and produce a prediction without access to future data,
often under strict latency constraints.

Most prior deep learning approaches to LOB prediction are \emph{episodic}---they
segment the data stream into fixed-length windows and apply batch-style
inference~\citep{zhang2019deeplob, wallbridge2020transformers}. While effective,
this paradigm introduces an artificial boundary between windows, discards
cross-window temporal dependencies, and is fundamentally misaligned with the
online nature of real market data. Recurrent architectures such as
LSTMs~\citep{hochreiter1997long} offer a closer fit to the streaming setting,
maintaining a hidden state updated at each tick, but they lack the capacity to
perform additional internal computation once a new observation has been
processed.

The Continuous Thought Machine (CTM)~\citep{sakana2025ctm}, introduced by
Sakana AI, proposes an alternative computational paradigm: rather than mapping
each input directly to an output in a single forward pass, the CTM allows the
model to perform a variable number of \emph{internal ticks} of computation
after receiving each external input. This decoupling of external observation
rate from internal processing rate is motivated by biological neural computation
and enables the model to reason more deeply before committing to a prediction.
The synchronization of neural activity histories across neuron pairs provides a
rich representational substrate that we hypothesize is particularly well-suited
to the non-stationary and event-driven dynamics of financial microstructure.

In this work, we evaluate whether the continuous-thought paradigm offers
measurable advantages for the LOB mid-price prediction task. We compare two
CTM-style architectures---a CTM-Inspired model that combines a GRU cell with
neural synchronization and internal pause ticks, and a full CTM implementation
with per-neuron MLPs and synchronization tensors---against logistic regression,
LSTM, and a Transformer window model. All models operate in an online,
sample-by-sample streaming setting on Bitcoin perpetual futures tick data from
Tardis.dev, resampled to 50\,ms intervals with a 5-second prediction horizon.

\paragraph{Contributions.}
\begin{itemize}
  \item We formulate LOB mid-price direction prediction as a strictly streaming
    classification task and provide a unified evaluation harness across
    episodic and recurrent model families.
  \item We introduce a CTM-Inspired architecture that adapts the internal-tick
    mechanism of the CTM to the streaming setting via a GRU backbone and
    stochastic neural synchronization pairs.
  \item We conduct controlled ablations over the number of internal ticks,
    synchronization pairs, prediction horizon, and resampling frequency,
    isolating the contribution of each CTM-specific component.
  \item We evaluate model robustness under synthetic market shocks---spread
    widening and sudden liquidity withdrawal---and characterize recovery time,
    prediction flip rate, and confidence deviation during the shock window.
\end{itemize}

The remainder of the paper is organized as follows. Section~\ref{sec:related}
reviews related work. Section~\ref{sec:data} describes the data pipeline and
feature construction. Section~\ref{sec:models} defines the model architectures.
Section~\ref{sec:experiments} presents experimental results and ablations.
Section~\ref{sec:shocks} reports shock robustness evaluations.
Section~\ref{sec:conclusion} concludes.

% ─────────────────────────────────────────────────────────────
\section{Related Work}
\label{sec:related}
% ─────────────────────────────────────────────────────────────

\subsection{Deep Learning for Limit Order Book Prediction}

The problem of mid-price prediction from LOB data has been approached with
increasingly expressive architectures. \citet{zhang2019deeplob} introduced
DeepLOB, a convolutional neural network that operates on fixed-length snapshots
of the full order book and demonstrated strong performance on the FI-2010
benchmark. Subsequent work explored recurrent variants to better capture
temporal dependencies~\citep{sirignano2019deep}, and Transformer-based models
have been applied to this setting to leverage long-range attention over LOB
sequences~\citep{wallbridge2020transformers}. A shared limitation of these
approaches is their reliance on fixed input windows, which introduces implicit
stationarity assumptions that may not hold during periods of market stress.

\subsection{Online and Streaming Learning}

Streaming learning---where a model processes one sample at a time and cannot
revisit past observations---has a long history in the machine learning
literature~\citep{bottou1998online}. In financial applications, online learning
is particularly relevant because market conditions shift non-stationarily and
batching introduces latency. Recurrent architectures are natural candidates for
this regime, as their hidden state serves as a compact sufficient statistic for
past observations. However, the depth of computation performed per step remains
fixed in standard RNN/LSTM architectures.

\subsection{Continuous Thought Machines}

The CTM~\citep{sakana2025ctm} introduces a neural architecture in which each
external observation triggers a configurable number of internal processing
steps before a prediction is emitted. Each neuron maintains a history of its
own pre-activations, and a synchronization readout computes pairwise dot
products over these histories to form the output. This design draws from
neuroscientific models of cortical computation where stimulus processing
involves recurrent activity sustained well beyond the initial sensory
input~\citep{lamme2000distinct}. The CTM has been evaluated on static
image classification and reasoning tasks; our work is the first to apply it
in a real-time financial time-series setting.

% ─────────────────────────────────────────────────────────────
\section{Data and Feature Construction}
\label{sec:data}
% ─────────────────────────────────────────────────────────────

\subsection{Raw Data}

We use Bitcoin perpetual futures (BTC-USDT) order book data sourced from
Tardis.dev. The raw data consists of incremental Level-2 order book updates,
recording each individual change to the book as a tuple
$(\text{timestamp}, \text{side}, \text{price}, \text{amount})$, where
timestamps are given in microseconds. We process these incremental updates
into fixed-interval snapshots using a custom orderbook rebuilder
(\texttt{tardis\_builder.py}) that maintains a sorted bid and ask dictionary
and emits a full $K$-level snapshot every $\Delta t = 50$\,ms.

\subsection{Resampling and Feature Extraction}

From the snapshot sequence, we compute a 45-dimensional feature vector at
each timestep. Let $b_k^t$ and $a_k^t$ denote the bid and ask prices at
depth level $k \in \{0, \ldots, K-1\}$ (with $K = 10$), and let
$\beta_k^t$, $\alpha_k^t$ denote the corresponding queue sizes. The mid-price
is $m^t = (b_0^t + a_0^t) / 2$ and the spread is
$s^t = a_0^t - b_0^t$. Features are:

\begin{align}
  \mathbf{x}^t = \bigl[
    &m^t,\; s^t,\; \log m^t,\; \log(m^t / m^{t-1}), \nonumber\\
    &\text{OI}^t,\;
     \{(m^t - b_k^t)/m^t\}_{k=0}^{K-1},\;
     \{(a_k^t - m^t)/m^t\}_{k=0}^{K-1}, \nonumber\\
    &\{\log(1 + \beta_k^t)\}_{k=0}^{K-1},\;
     \{\log(1 + \alpha_k^t)\}_{k=0}^{K-1}
  \bigr]
\end{align}

where the order imbalance is
\[
  \text{OI}^t = \frac{\sum_k \beta_k^t - \sum_k \alpha_k^t}
                     {\sum_k \beta_k^t + \sum_k \alpha_k^t + \epsilon}.
\]

This yields $5 + 4K = 45$ features per timestep.

\subsection{Label Construction}

We assign a binary direction label to each timestep based on the mid-price
$H = 5$\,s in the future:
\[
  \delta^t = \frac{m^{t+H} - m^t}{m^t}, \qquad
  y^t = \mathbf{1}[\delta^t > 0].
\]

To avoid polluting the class distribution with effectively flat samples, we
apply an $\epsilon$-filter and only retain timesteps for which
$|\delta^t| > \epsilon$, with $\epsilon = 10^{-4}$ (a relative threshold of
one basis point). The resulting dataset is approximately balanced between up
and down moves.

% ─────────────────────────────────────────────────────────────
\section{Model Architectures}
\label{sec:models}
% ─────────────────────────────────────────────────────────────

All models receive input $\mathbf{x}^t \in \mathbb{R}^{45}$ at each timestep
and produce a scalar logit $z^t \in \mathbb{R}$, from which the predicted
probability is $\hat{p}^t = \sigma(z^t)$. The binary prediction is
$\hat{y}^t = \mathbf{1}[\hat{p}^t > 0.5]$.

\subsection{Logistic Baseline}

The logistic baseline is a stateless two-layer MLP:
\[
  z^t = W_2 \,\text{ReLU}(W_1 \mathbf{x}^t + b_1) + b_2,
\]
with $W_1 \in \mathbb{R}^{64 \times 45}$ and $W_2 \in \mathbb{R}^{1 \times 64}$.
Each timestep is processed independently with no memory of prior observations.

\subsection{LSTM}

The LSTM baseline~\citep{hochreiter1997long} maintains a hidden state
$(h^t, c^t) \in \mathbb{R}^{64} \times \mathbb{R}^{64}$ updated at each step:
\[
  (h^t, c^t) = \text{LSTM}(\mathbf{x}^t,\; h^{t-1},\; c^{t-1}),
  \qquad z^t = W_\text{out} h^t.
\]

\subsection{Transformer}

The Transformer baseline operates on a sliding window of length $W = 40$
timesteps. Input features are linearly projected to a $d_\text{model} = 64$
dimensional embedding space and summed with learned positional embeddings.
A two-layer Transformer encoder with $n_\text{head} = 4$ heads and a causal
attention mask is applied, and the representation at the final position is
passed through a linear output head. This is the only model in our suite
that is episodic rather than streaming.

\subsection{CTM-Inspired}

Our CTM-Inspired model adapts the key mechanisms of the CTM to a streaming
GRU backbone. It maintains three components of state: a GRU hidden vector
$h^t \in \mathbb{R}^{64}$, a circular activity buffer
$B^t \in \mathbb{R}^{L \times 64}$ (with $L = 64$), and a buffer write
pointer $p^t$.

Upon receiving $\mathbf{x}^t$, the model first performs one external update:
\[
  h^t_1 = \text{GRUCell}(\mathbf{x}^t, h^{t-1}),
\]
writes $h^t_1$ into $B^t$, and then performs $T - 1$ additional
\emph{internal ticks} using a learned pause embedding
$\mathbf{e}_\text{pause} \in \mathbb{R}^{45}$ in place of the external input:
\[
  h^t_\tau = \text{GRUCell}(\mathbf{e}_\text{pause},\; h^t_{\tau-1}),
  \quad \tau = 2, \ldots, T,
\]
each writing its activation into the buffer. The output logit at the final
tick is derived from a \emph{neural synchronization} readout: $P = 128$
pairs of buffer positions $(i_j, j_j)$ are sampled uniformly, their
representations are dot-producted, and the resulting synchronization vector
$\mathbf{s}^t \in \mathbb{R}^{P}$ is passed through a small MLP:
\[
  s^t_j = \langle B^t[i_j],\; B^t[j_j] \rangle, \quad
  z^t = \text{MLP}_\text{sync}(\mathbf{s}^t).
\]
By default $T = 5$ internal ticks are used.

\subsection{CTM (Full)}

The full CTM implementation replaces the GRU cell with a dedicated synapse
MLP and assigns each of $D = 64$ neurons its own activation function
parameterized by a small per-neuron MLP operating on the neuron's pre-activation
history $\mathbf{h}^t_d \in \mathbb{R}^{m}$ (with history length $m = 8$):
\[
  \tilde{a}^t = \text{SynapseMLP}([\mathbf{a}^{t-1};\, \mathbf{x}^t]),
  \quad a^t_d = \text{NeuronMLP}_d(\mathbf{h}^t_d),
\]
where $\mathbf{a}^t \in \mathbb{R}^{D}$ is the post-activation vector and
the history buffer is updated by appending $\tilde{a}^t_d$. The output
uses the same stochastic synchronization readout as the CTM-Inspired model,
with $T = 10$ internal ticks.

% ─────────────────────────────────────────────────────────────
\section{Experiments}
\label{sec:experiments}
% ─────────────────────────────────────────────────────────────

\subsection{Setup}

All models are trained online with batch size 1 using the Adam
optimizer~\citep{kingma2015adam} with learning rate $10^{-4}$ and weight
decay $10^{-5}$, minimizing binary cross-entropy. Training proceeds
chronologically for up to 20 epochs with early stopping (patience 10) based
on validation accuracy. The dataset is split 80/20 chronologically, with the
first 80\% used for training and the remaining 20\% for validation and
evaluation.

\subsection{Evaluation Metrics}

We report three metrics on the held-out validation set:
\begin{itemize}
  \item \textbf{Accuracy}: fraction of correctly classified direction labels.
  \item \textbf{AUROC}: area under the receiver operating characteristic
    curve, measuring ranking quality independent of the decision threshold.
  \item \textbf{Flip Rate}: the fraction of consecutive prediction pairs
    $({\hat{y}^t, \hat{y}^{t+1}})$ that differ, measuring prediction stability.
    A high flip rate indicates erratic behavior unsuitable for practical use.
\end{itemize}

\subsection{Main Results}

Table~\ref{tab:main_results} summarizes model performance. Results are
reported on the held-out 20\% chronological split.

\begin{table}[h]
\centering
\caption{Main results on held-out test set ($H = 5$\,s, $\Delta t = 50$\,ms,
         $\epsilon = 10^{-4}$).}
\label{tab:main_results}
\begin{tabular}{lccc}
  \toprule
  Model              & Accuracy ($\uparrow$) & AUROC ($\uparrow$) & Flip Rate ($\downarrow$) \\
  \midrule
  Logistic           & ---                   & ---                & ---                      \\
  LSTM               & ---                   & ---                & ---                      \\
  Transformer        & ---                   & ---                & ---                      \\
  CTM-Inspired       & ---                   & ---                & ---                      \\
  CTM (Full)         & ---                   & ---                & ---                      \\
  \bottomrule
\end{tabular}
\end{table}

\subsection{Ablations}

\paragraph{Internal Ticks.}
We ablate the number of internal ticks $T \in \{1, 2, 5, 10\}$ for the
CTM-Inspired model, holding all other hyperparameters fixed. Tick $T = 1$
reduces the model to a standard GRU with a synchronization readout, while
higher values allow progressively more internal deliberation per observation.

\paragraph{Synchronization Pairs.}
We vary the number of synchronization pairs $P \in \{32, 64, 128, 256\}$,
controlling the richness of the activity-history readout. Too few pairs
under-sample the buffer; too many may introduce variance.

\paragraph{Resampling Frequency.}
We compare $\Delta t \in \{10, 50\}$\,ms to assess sensitivity to the
input granularity. Finer resolution exposes the model to more raw book
dynamics but increases the sequence length proportionally.

\paragraph{Prediction Horizon.}
We vary $H \in \{0.5, 1.0, 5.0\}$\,s. Shorter horizons make prediction
harder due to greater noise, while longer horizons are more forecastable
but of less practical relevance for high-frequency strategies.

\begin{table}[h]
\centering
\caption{Ablation study on CTM-Inspired (accuracy on held-out set).}
\label{tab:ablations}
\begin{tabular}{llc}
  \toprule
  Parameter              & Value       & Accuracy ($\uparrow$) \\
  \midrule
  \multirow{4}{*}{Ticks $T$}
                         & 1           & ---                   \\
                         & 2           & ---                   \\
                         & 5 (default) & ---                   \\
                         & 10          & ---                   \\
  \midrule
  \multirow{4}{*}{Sync Pairs $P$}
                         & 32          & ---                   \\
                         & 64          & ---                   \\
                         & 128 (default) & ---                 \\
                         & 256         & ---                   \\
  \midrule
  \multirow{2}{*}{$\Delta t$}
                         & 10\,ms      & ---                   \\
                         & 50\,ms (default) & ---              \\
  \midrule
  \multirow{3}{*}{Horizon $H$}
                         & 0.5\,s      & ---                   \\
                         & 1.0\,s      & ---                   \\
                         & 5.0\,s (default) & ---              \\
  \bottomrule
\end{tabular}
\end{table}

% ─────────────────────────────────────────────────────────────
\section{Shock Robustness Evaluation}
\label{sec:shocks}
% ─────────────────────────────────────────────────────────────

\subsection{Shock Types}

Market microstructure can experience abrupt structural changes that fall
outside the distribution of normal trading. To evaluate model robustness, we
inject two types of synthetic shocks into a held-out clean data segment and
compare model behavior before, during, and after the shock window.

\paragraph{Spread Widening.}
The bid-ask spread is artificially widened by a factor of $3\times$ for a
contiguous window of 40 timesteps (2\,s at 50\,ms resolution) centered at
the midpoint of the evaluation segment. The shock is applied symmetrically:
the ask side is moved up and the bid side is moved down by equal amounts,
preserving the mid-price.

\paragraph{Liquidity Cliff.}
The top three levels of the order book on both the bid and ask sides are
zeroed out for the same 40-step window, simulating a sudden withdrawal of
resting liquidity. This shock leaves prices unchanged but dramatically
alters the size and imbalance features that the models rely on.

\subsection{Evaluation Protocol}

Pre-trained models are evaluated in two conditions: (i) on the clean segment
without any modification, and (ii) on the same segment with the synthetic
shock injected. We report three shock-specific metrics:
\begin{itemize}
  \item \textbf{Accuracy during shock}: classification accuracy restricted to
    timesteps within the shock window.
  \item \textbf{Flip rate during shock}: prediction instability within the
    shock window, which should ideally remain low.
  \item \textbf{Recovery time}: the number of timesteps after the shock ends
    before the model's predictions return to their pre-shock behavior, measured
    by the rolling accuracy returning within 1\% of its pre-shock value.
\end{itemize}

\subsection{Results}

Table~\ref{tab:shock_results} summarizes model robustness under each shock type.

\begin{table}[h]
\centering
\caption{Shock robustness evaluation. Acc.\,Shock = accuracy during shock window;
         Flip\,Shock = flip rate during shock window;
         Recovery = timesteps to recover post-shock.}
\label{tab:shock_results}
\begin{tabular}{lcccccc}
  \toprule
  & \multicolumn{3}{c}{Spread Widening} & \multicolumn{3}{c}{Liquidity Cliff} \\
  \cmidrule(lr){2-4} \cmidrule(lr){5-7}
  Model         & Acc. & Flip & Rec. & Acc. & Flip & Rec. \\
  \midrule
  Logistic      & ---  & ---  & ---  & ---  & ---  & ---  \\
  LSTM          & ---  & ---  & ---  & ---  & ---  & ---  \\
  Transformer   & ---  & ---  & ---  & ---  & ---  & ---  \\
  CTM-Inspired  & ---  & ---  & ---  & ---  & ---  & ---  \\
  CTM (Full)    & ---  & ---  & ---  & ---  & ---  & ---  \\
  \bottomrule
\end{tabular}
\end{table}

% ─────────────────────────────────────────────────────────────
\section{Conclusion}
\label{sec:conclusion}
% ─────────────────────────────────────────────────────────────

We have presented a systematic evaluation of the Continuous Thought Machine
paradigm in the context of streaming limit order book prediction. By formulating
the task as a strictly online classification problem and benchmarking CTM-style
architectures against logistic regression, LSTM, and Transformer baselines on
high-resolution Bitcoin perpetual futures data from Tardis.dev, we isolate the
contribution of internal ticking and neural synchronization to predictive
performance and stability.

Our CTM-Inspired architecture demonstrates that the core mechanisms of the
CTM---internal pause ticks and pairwise synchronization of activity histories---
can be grafted onto a recurrent backbone with minimal overhead, and that the
number of internal ticks $T$ can be tuned as a hyperparameter to trade
inference latency for representational depth. The shock robustness evaluation
provides a complementary view of model reliability, quantifying how quickly
each architecture recovers from out-of-distribution perturbations to the
order book structure.

\paragraph{Limitations.}
The current evaluation is limited to a single asset and trading venue.
Generalization to other instruments, venues, and market regimes remains to be
studied. The full CTM implementation uses stochastic synchronization pair
sampling, which introduces variance in both training and evaluation. Future work
should explore deterministic synchronization readouts and learned pair
selection.

\paragraph{Future Work.}
Promising directions include applying the CTM to tick-level event streams
rather than fixed-interval snapshots, integrating the CTM's confidence signal
(derived from synchronization entropy) as a live uncertainty estimate for
downstream risk management, and extending the architecture to multi-asset
prediction with shared internal representations.

% ─────────────────────────────────────────────────────────────
\bibliographystyle{plainnat}
\bibliography{references}
% ─────────────────────────────────────────────────────────────

\end{document}
